% !TEX program = xelatex
% ============================================================
%  ورقة بيانات — Datasheet Template
%  قالب لورقة بيانات المنتج التقنية (Datasheet) بمظهر تقني مختصر.
%  Compile with: xelatex main.tex (run twice)
% ============================================================
%  Features:
%    - Product name header + model placeholder
%    - Key features bullet list
%    - Specifications table (Parameter | Value | Unit | Notes)
%    - Typical application diagram (TikZ placeholder)
%    - Mechanical drawing placeholder
%    - Ordering information table, contact footer
% ============================================================

\documentclass[11pt,a4paper]{article}

\usepackage[a4paper,margin=2cm,top=2.2cm,bottom=2.2cm,headheight=16pt]{geometry}
\usepackage{graphicx}
\usepackage{array}
\usepackage{booktabs}
\usepackage{tabularx}
\usepackage{multirow}
\usepackage{enumitem}
\usepackage{float}
\usepackage{caption}
\usepackage{titlesec}
\usepackage{fancyhdr}
\usepackage{setspace}
\usepackage[table]{xcolor}
\usepackage{amsmath}
\usepackage{amssymb}
\usepackage{tikz}
\usetikzlibrary{positioning,calc,shapes.geometric,arrows.meta}

\usepackage{bidi}
\usepackage[no-math]{fontspec}
\usepackage[quiet]{polyglossia}

\setmainfont[Scale=1]{NotoKufiArabic-Regular.ttf}
\newfontfamily\arabicfont[Script=Arabic,Scale=0.95]{NotoKufiArabic-Regular.ttf}
\newfontfamily\englishfont[Scale=0.95]{TeX Gyre Termes}

\setdefaultlanguage[calendar=gregorian,hijricorrection=1]{arabic}
\setotherlanguage[variant=british]{english}

\usepackage[breaklinks,hidelinks]{hyperref}
\hypersetup{pdftitle={ورقة بيانات},pdfauthor={اسم الشركة}}

\addto\captionsarabic{%
  \renewcommand{\contentsname}{الفهرس}%
  \renewcommand{\figurename}{الشكل}%
  \renewcommand{\tablename}{الجدول}%
}

\titleformat{\section}{\large\bfseries}{\thesection}{0.7em}{}
\titlespacing*{\section}{0pt}{1.0ex plus 0.3ex}{0.6ex plus 0.2ex}

\pagestyle{fancy}
\fancyhf{}
\fancyhead[R]{\small\itshape ورقة بيانات (Datasheet)}
\fancyhead[L]{\small اسم المنتج --- \LR{Model-XXX}}
\fancyfoot[C]{\small\thepage}
\renewcommand{\headrulewidth}{0.3pt}

\setlist[itemize]{topsep=0.2em, itemsep=0.1em, parsep=0pt, leftmargin=1.4em}
\setlist[enumerate]{topsep=0.2em, itemsep=0.1em, parsep=0pt, leftmargin=1.4em}

\newcolumntype{L}[1]{>{\raggedright\arraybackslash}p{#1}}
\newcolumntype{C}[1]{>{\centering\arraybackslash}p{#1}}
\newcolumntype{R}[1]{>{\raggedleft\arraybackslash}p{#1}}

\definecolor{brandColor}{HTML}{1A3C6B}
\definecolor{headerBg}{HTML}{1A3C6B}
\definecolor{altRow}{HTML}{F0F4F8}
\definecolor{blockFill}{HTML}{E8EEF4}
\definecolor{accentColor}{HTML}{2E6DA4}

\setstretch{1.2}

\begin{document}

% ============================================================
% PRODUCT HEADER
% ============================================================
\noindent
\begin{tikzpicture}
\fill[headerBg] (0,0) rectangle (\textwidth,1.6cm);
\node[anchor=west, text=white, font=\LARGE\bfseries] at (0.4cm,1.0cm) {اسم المنتج};
\node[anchor=west, text=white, font=\normalsize] at (0.4cm,0.35cm) {ورقة بيانات المنتج (Product Datasheet)};
\node[anchor=east, text=white, font=\Large\bfseries] at (\textwidth-0.4cm,0.7cm) {\LR{Model-XXX}};
\end{tikzpicture}

\vspace{0.6em}

\noindent\textbf{اسم الشركة:} --- \hfill \textbf{الإصدار:} \LR{v1.0} \hfill \textbf{التاريخ:} ---

\vspace{0.8em}

% ============================================================
% 1. PRODUCT DESCRIPTION
% ============================================================
\section{وصف المنتج (Product Description)}
% TODO: اكتب وصفاً مختصراً للمنتج.
وصف مختصر للمنتج ووظيفته الرئيسية والاستخدام المقصود. يبين هذا القسم الغرض
العام للمنتج والفئة التي ينتمي إليها.

% ============================================================
% 2. KEY FEATURES
% ============================================================
\section{الميزات الرئيسية (Key Features)}
% TODO: اذكر الميزات الرئيسية للمنتج.
\begin{itemize}
    \item ميزة رئيسية أولى مع وصف مختصر.
    \item ميزة رئيسية ثانية مع وصف مختصر.
    \item ميزة رئيسية ثالثة مع وصف مختصر.
    \item ميزة رئيسية رابعة مع وصف مختصر.
    \item ميزة رئيسية خامسة مع وصف مختصر.
\end{itemize}

% ============================================================
% 3. SPECIFICATIONS
% ============================================================
\section{المواصفات (Specifications)}
% TODO: املأ جدول المواصفات بالقيم الفعلية.
\rowcolors{2}{altRow}{white}
\begin{table}[H]
\centering
\renewcommand{\arraystretch}{1.25}
\begin{tabularx}{\textwidth}{L{4.5cm} L{3.5cm} C{2cm} X}
\toprule
\textbf{المعامل (Parameter)} & \textbf{القيمة (Value)} & \textbf{الوحدة (Unit)} & \textbf{ملاحظات (Notes)} \\
\midrule
جهد التشغيل (\LR{Supply Voltage}) & \LR{3.3 -- 5.0}  & \LR{V}   & نطاق الجهد المسموح. \\
تيار الاستهلاك (\LR{Current Consumption}) & \LR{120} & \LR{mA}  & في الوضع النشط. \\
تيار السكون (\LR{Standby Current}) & \LR{5}    & \LR{\textmu A} & في وضع السكون. \\
درجة حرارة التشغيل (\LR{Operating Temp.}) & \LR{-40 -- +85} & \LR{\textdegree C} & نطاق الحرارة. \\
التردد (\LR{Clock Frequency}) & \LR{16}   & \LR{MHz}  & تردد الساعة الداخلي. \\
الذاكرة (\LR{Flash Memory}) & \LR{512}  & \LR{KB}   & ذاكرة البرنامج. \\
الذاكرة العشوائية (\LR{RAM}) & \LR{64}   & \LR{KB}   & ذاكرة البيانات. \\
الأبعاد (\LR{Dimensions}) & \LR{48 \texttimes 25} & \LR{mm}  & الطول \texttimes العرض. \\
الوزن (\LR{Weight}) & \LR{12}   & \LR{g}    & بدون التغليف. \\
الواجهة (\LR{Interface}) & \LR{UART / I2C / SPI} & --- & الواجهات المدعومة. \\
\bottomrule
\end{tabularx}
\end{table}
\rowcolors{1}{}{}

% ============================================================
% 4. TYPICAL APPLICATION DIAGRAM
% ============================================================
\section{مخطط التطبيق النموذجي (Typical Application)}
% TODO: استبدل المخطط المؤقت بمخطط التطبيق الفعلي.
\begin{figure}[H]
\centering
\begin{tikzpicture}[
  every node/.style={font=\small},
  block/.style={fill=blockFill,draw=accentColor,thick,rounded corners=2pt,inner sep=5pt,minimum width=2.4cm,minimum height=1.1cm,align=center},
  arrow/.style={draw=accentColor,thick,-{Stealth[length=3.5mm]}},
]
\node[block] (src)   at (0,0)   {المصدر\\(\LR{Source})};
\node[block] (mcu)   at (4,0)   {المنتج\\(\LR{Model-XXX})};
\node[block] (load)  at (8,0)   {الحمل\\(\LR{Load})};
\node[block] (sense) at (4,-2.2){الحساس\\(\LR{Sensor})};
\draw[arrow] (src)  -- (mcu);
\draw[arrow] (mcu)  -- (load);
\draw[arrow] (sense) -- (mcu);
\end{tikzpicture}
\caption{مخطط التطبيق النموذجي (Typical Application Diagram)}
\end{figure}

% ============================================================
% 5. MECHANICAL DRAWING
% ============================================================
\section{الرسم الميكانيكي (Mechanical Drawing)}
% TODO: أدرج الرسم الميكانيكي الفعلي للمنتج.
\begin{figure}[H]
\centering
\begin{tikzpicture}
\draw[thick,accentColor] (0,0) rectangle (10,4);
\draw[gray!50,dashed] (0,0) grid[step=0.5] (10,4);
\node[font=\large\color{accentColor}] at (5,2) {موضع الرسم الميكانيكي (\LR{Mechanical Drawing})};
\node[font=\small\color{gray}] at (5,1.4) {\LR{48 mm \texttimes\ 25 mm}};
\end{tikzpicture}
\caption{الرسم الميكانيكي للأبعاد (Mechanical Drawing Placeholder)}
\end{figure}

% ============================================================
% 6. ORDERING INFORMATION
% ============================================================
\section{معلومات الطلب (Ordering Information)}
% TODO: املأ جدول معلومات الطلب.
\begin{table}[H]
\centering
\renewcommand{\arraystretch}{1.3}
\begin{tabularx}{\textwidth}{L{4cm} L{4cm} X}
\toprule
\textbf{رقم الجزء (Part Number)} & \textbf{الوصف (Description)} & \textbf{التغليف (Package)} \\
\midrule
\LR{Model-XXX-A} & النموذج الأساسي & \LR{Tray} \\
\LR{Model-XXX-B} & النموذج الموسع  & \LR{Tape \& Reel} \\
\LR{Model-XXX-C} & النموذج الصناعي & \LR{Tray} \\
\bottomrule
\end{tabularx}
\end{table}

% ============================================================
% CONTACT FOOTER
% ============================================================
\vspace{0.8em}
\noindent
\begin{tikzpicture}
\fill[headerBg] (0,0) rectangle (\textwidth,1.3cm);
\node[anchor=west, text=white, font=\small\bfseries] at (0.4cm,0.95cm) {معلومات الاتصال (Contact Information)};
\node[anchor=west, text=white, font=\footnotesize] at (0.4cm,0.4cm) {اسم الشركة --- \LR{www.example.com} --- \LR{support@example.com}};
\node[anchor=east, text=white, font=\footnotesize] at (\textwidth-0.4cm,0.4cm) {\LR{\textcopyright\ 2025}};
\end{tikzpicture}

\end{document}
